\chapter{Using \meq}
\label{ch:using}

{\meq} is a library first and a program second. This chapter is about the
program: what it needs in order to build, how a run is described, what it
writes, and what to do when a solve does not converge. Calling the library
directly is \chref{ch:embedding}.

\section{Building}

\subsection{The one hard dependency}

{\meq} is C++17 and CMake, and it has a single dependency that is not
straightforward: MFEM \cite{Anderson2021}, the finite element library the
discretisation is built on.

\begin{quote}
	\textbf{A release of MFEM will not do, and neither will \texttt{master}.}
\end{quote}

{\meq} needs the mixed \texttt{DarcyForm} and the hybridizable discontinuous
Galerkin integrators that go with it, and for the curved boundary of
\S\ref{sec:curved} it needs the transfer-path machinery that implements
Cockburn--Solano extension \cite{CockburnSolano2012}. That work lives on
development branches, and \texttt{INSTALL.md} in the source tree names them and
gives the merge recipe. This is the one place where building {\meq} is more
than a package manager away, and there is no way round it: the method is the
branch.

Everything else is ordinary. A recent compiler; CMake; the header-only TOML
parser toml11, vendored as a git submodule so that its version is {\meq}'s to
control rather than the machine's; Boost and Eigen, both header-only and both
needed by the library rather than only by the tests --- Boost.Math supplies the
Jacobi polynomials the Zernike basis of \chref{ch:embedding} is evaluated
through, and Eigen supplies the singular value decomposition its fits are solved
with. NetCDF's C++ bindings are optional; without them the library is built
without its gridded writer and the program cannot produce its interchange
output. Boost.Test is needed only to build the tests, and \texttt{clang-tidy}
only to run the naming check.

Several libraries that a reader might expect to configure --- SuiteSparse,
SUNDIALS, GSLIB, LAPACK, MKL, CUDA --- are reached \emph{through} MFEM rather
than found separately. {\meq} reads which of them are present out of MFEM's own
configuration file and adapts, so the configure summary reports what MFEM was
built with rather than what {\meq} guessed.

\subsection{Configure and build}

\begin{verbatim}
    git submodule update --init --recursive
    cmake -B build
    cmake --build build -j4
    cd build && ctest --output-on-failure
\end{verbatim}

\texttt{MFEM\_DIR} is how the library is found. It is a cache variable, it also
reads the environment, and it has a default, so on a machine laid out the
obvious way no \texttt{-D} is needed at all. The finder accepts either layout an
MFEM build leaves behind --- an installed tree, or an in-source \texttt{make}
build --- and needs no flag to tell them apart.

\begin{quote}
	\textbf{Never run a bare \texttt{make -j} or \texttt{cmake -{}-build -j}
	with no number.} With no argument the job count is unbounded and goes to
	the host's core count, which on a container or a virtual machine is not the
	same thing as the memory behind it. Give a number.
\end{quote}

Three configure-time options are worth knowing. \texttt{MEQ\_BUILD\_APP}, on by
default, builds the program; \texttt{MEQ\_BUILD\_TESTS}, likewise, builds the
test binaries and registers them; and \texttt{MEQ\_ENABLE\_COVERAGE}
instruments the build. That last is an option rather than a build type on
purpose: a coverage build wants no optimisation and a release build wants all of
it, and the difference should be chosen rather than inherited from whatever
\texttt{CMAKE\_BUILD\_TYPE} happens to hold.

A debug build of {\meq} buys the debugger and nothing else. It is worth saying
so, because the opposite is a natural assumption. MFEM's internal consistency
checks are gated on its own debug flag rather than on \texttt{NDEBUG}, and a
normal MFEM install has that flag off --- so those assertions are already dead
in every {\meq} build, and no flag passed to {\meq} revives one. Recovering them
means building a second MFEM configured for debugging and pointing at it, which
is a real option when chasing a suspected contract violation with the library
and a much larger undertaking than a build type.

\subsection{Building without MFEM}

MFEM is deliberately not a required package. Without it, CMake builds the
MFEM-free half of the library --- the configuration parser, the profiles, the
source terms, the boundary shapes --- skips every file that includes the
library, and skips every test that needs a solver, saying so for each. That is
not a convenience: it is what makes continuous integration possible at all,
since the branch {\meq} needs is not published anywhere a hosted runner can
fetch it. The corollary is that a green badge is evidence about the
configuration and profile layers and about nothing else.

\section{Running}

\begin{verbatim}
    meq config.toml     solve the equilibrium
    meq --help          usage
    meq --version       the build this is
\end{verbatim}

There are no other flags and no subcommands, and that is a decision rather than
an unfinished corner. \emph{A configuration file is a record of what was run; a
command line is not.} When somebody asks six months later which equilibrium a
plot came from, the answer should be a file that can be re-run --- and every
output {\meq} writes carries the name of the file that produced it.

\subsection{What a run does}

In order:

\begin{enumerate}

\item \textbf{Reads the configuration}, and validates all of it before doing any
	work. An unknown key is an error, not something ignored, and so is an
	unknown table; the message names the nearest accepted spelling. A key that
	is silently ignored is a run that quietly did something other than what the
	file says, which is the failure this behaviour exists to prevent.

\item \textbf{Builds the mesh}, either by triangulating the box given in the
	configuration or by reading a mesh file, and applies the requested uniform
	refinements.

\item \textbf{Builds the source term} $F$, and with it $\partial F/\partial\psi$.

\item \textbf{Builds the computational domain.} With no shape given, this is the
	mesh itself and $\psi = 0$ is imposed on its boundary: the \emph{fitted}
	path. With a shape, {\meq} marks the background elements lying inside
	$\Gamma$, extracts them as a submesh, and builds the transfer paths that
	carry the boundary condition outward: the \emph{curved} path of
	\S\ref{sec:curved}.

\item \textbf{Solves}, by Newton --- possibly inside an adaptive loop, and
	possibly with the flux on axis as an additional unknown.

\item \textbf{Post-processes}, building the enriched potential $\psi^{\star}$.

\item \textbf{Writes the answer}, three times.

\end{enumerate}

\subsection{What it prints, and what to read}

Progress goes to standard output and diagnostics to standard error, every line
prefixed so that a log can be filtered. The line that matters most on an
ordinary run says how many Newton iterations it took, on how many elements, at
what degree.

Under it is the residual history, printed on \emph{every} run and not only on
failure, with a column giving the observed order of convergence. That column is
there for one reason, and it is the reason given in \chref{ch:intro}: a Jacobian
with an error in it still converges to the right answer at the right rate in
$h$, and the only thing that changes is the path. The observed order falling
from two to one is the single diagnostic that sees it. If you have written your
own source term --- \chref{ch:embedding} --- this column is what tells you
whether its derivative is right.

An adaptive run additionally prints one row per cycle, and says which stopping
condition ended it. A run with the axis flux as an unknown prints the converged
value and the residual of the constraint that defines it.

\subsection{Exit codes}

\begin{center}
\begin{tabular}{c l p{0.55\textwidth}}
	\hline
	Code & Name & Meaning \\
	\hline
	0 & solved & The equilibrium was computed and everything asked for was
		written. \\
	1 & configuration error & Nothing was attempted. The file was unreadable or
		invalid, a key was wrong, the geometry could not be built, or the file
		asked for something {\meq} refuses. \\
	2 & solve failed & The configuration was good and the nonlinear iteration
		did not converge, including after the fallback. The residual history is
		printed, so the failure is diagnosable rather than merely reported. \\
	3 & output failed & The equilibrium was computed and something went wrong
		writing it --- most often an output directory that does not exist.
		{\meq} does not create the output directory. \\
	\hline
\end{tabular}
\end{center}

The split between~1 and~2 is the useful one for a script. A~1 means the input is
wrong and re-running will not help; a~2 means the input is well formed and the
solver could not do it, which is a case for refining the mesh, raising the
degree or supplying an initial guess.

There is exactly one thing {\meq} refuses rather than approximating: a boundary
condition given as the source's exact solution. That is a convergence-study
device, and a \texttt{Source} carries $F$ and $\partial F/\partial\psi$ and no
closed form --- the closed forms live in the test fixtures, where they are the
thing being converged against. The program says so and exits~1; the library
still offers it.

\section{The configuration file}

A run is described entirely by a TOML file. Three tables are required and the
rest are optional, and an absent optional table behaves exactly like a present
but empty one. Appendix~\ref{ap:config} documents every table and every key.

\begin{verbatim}
    [mesh]
    RMin = 0.6
    RMax = 1.4
    ZMin = -0.6
    ZMax = 0.6
    NR = 8
    NZ = 8

    [discretisation]
    PolynomialDegree = 3

    [source]
    Type = "soloviev"
    A = -0.52
\end{verbatim}

Two conventions are worth stating here because they surprise people. Table names
are lower case and \emph{keys are \texttt{UpperCamelCase}}, which is a
deliberate mismatch with the C++ naming rule used everywhere else in the tree
and is not an oversight to be reconciled. And numbers accept either TOML
spelling --- \texttt{RMin = 0} and \texttt{RMin = 0.0} both work --- while
counts do not, so \texttt{NR = 4.0} is refused rather than truncated. That last
distinction is worth more than it looks: the natural way to accept both
spellings in the TOML library used here returns the \emph{default} when a
conversion fails, because a failed conversion is indistinguishable from a
missing key, which turns a loud failure into a silent wrong answer in the
configuration --- the last place anybody wants one.

\section{Choosing the discretisation}

Two numbers set the accuracy: the polynomial degree $k$ and the mesh spacing
$h$. Both $\psi$ and the flux $q$ converge at $k+1$, and the enriched potential
$\psi^{\star}$ at $k+2$.

Raising $k$ is usually the better buy, for the reason \chref{ch:intro} gives:
hybridization leaves one globally coupled system on the mesh \emph{faces}, so
the cost of the global solve grows with the face count rather than with the
volume unknowns, and the extra volume unknowns a higher degree brings are
eliminated element by element. Refining the mesh and raising the degree are also
independent cures for a solve that will not converge, and \S\ref{sec:notconverge}
returns to that.

There are two ceilings. Beyond about $k = 5$, or eight refinements, round-off
dominates and a convergence table that flattens there is behaving correctly
rather than revealing a bug. And a study on a domain with a re-entrant or merely
non-smooth boundary is capped by the geometry rather than by the elements: on a
polygon approximating a smooth curve, the corner singularity at each vertex caps
the achievable order well below the design order whatever $k$ is, which is
exactly the difficulty \S\ref{sec:curved} exists to remove.

The stabilisation $\tau$ is the third parameter and its default is~1. Both
source papers set it there and note that optimal order needs only
$\tau = O(1)$.

\begin{quote}
	\textbf{Keep $\tau$ constant.} A solution-dependent stabilisation must also
	supply its own derivative for the Newton Jacobian, and omitting that gives
	no wrong answer --- only slow Newton convergence, which is a failure that
	survives a passing regression suite. A constant $\tau$ never reaches that
	code path. Deriving $\tau$ from the local diffusion coefficient is
	specifically not advised.
\end{quote}

\section{The plasma boundary}
\label{sec:curved}

The plasma boundary $\Gamma$ is a smooth curve. A mesh is not. {\meq} offers
two ways of reconciling them, and only one of them is any good.

The \emph{fitted} path takes $\Gamma$ to be the mesh boundary and imposes
$\psi = 0$ on it. It is the right thing for a rectangular benchmark domain and
for convergence studies against a closed form. It is the \emph{wrong} thing for
a real plasma shape, because a polygonal approximation to a curved boundary caps
the convergence rate whatever the polynomial degree: each vertex carries a
corner singularity, interior pollution spreads it, and the measured order on a
fine polygon at high degree comes out at roughly half the design order. No
accuracy inside the elements recovers that, because there is nothing wrong with
the elements --- the \emph{domain} is wrong.

The \emph{curved} path does not mesh the boundary at all. Give the run a shape
and {\meq} marks the background elements lying entirely inside $\Gamma$,
extracts them as a computational domain $\Omega_h$ whose boundary $\Gamma_h$ is
therefore inscribed, builds a short path from each point of $\Gamma_h$ out to
$\Gamma$, and imposes on $\Gamma_h$ not the boundary condition itself but the
boundary condition carried back along the path. The line integral that carries
it uses the computed flux $q$, so it is available at the flux's own order: the
mixed method paying off again. The property that makes this more than a
heuristic is that optimal order survives a gap of only $O(h)$
\cite{CockburnSolano2012}.

Two shape families are built in: the Miller parametrisation
\cite{Miller1998}, in which a plasma cross-section is described by a major
radius, a minor radius, an elongation, a triangularity and a squareness; and its
extended Miller harmonic generalisation \cite{ArbonCandyBelli2021}, which
replaces the two shaping parameters by Fourier coefficients and can represent
shapes Miller cannot.

\begin{quote}
	\textbf{The curved path leaves a band between $\Gamma_h$ and $\Gamma$ that
	is inside the plasma and outside the mesh}, $O(h)$ wide. It is real and it
	is visible in the output; \S\ref{sec:output} says what each format does
	about it.
\end{quote}

There is one further consequence, which matters only if you go looking. On the
fitted path the trace matrix is symmetric to round-off and negative definite. On
the curved path it is \emph{not symmetric}, and that is not a defect --- it is
what the transfer technique is. The continuous Grad--Shafranov operator is
self-adjoint; the same operator with a \emph{transferred} Dirichlet datum is
not, because the datum on a face depends on the flux along a path leaving the
element. So {\meq}'s headline configuration is genuinely non-symmetric and an
unsymmetric sparse LU is the right solver for it.

\section{Adaptive refinement}

Given a shape or not, {\meq} will run the usual loop --- solve, post-process,
estimate, mark, refine --- against the residual error estimator of
S\'anchez-Vizuet, Solano and Cerfon \cite{SanchezVizuet2020adaptive}, stopping at
a target error or an iteration limit and saying which. Two marking strategies
are offered, D\"orfler's bulk-chasing rule and a simple fraction-of-maximum rule.
They read their one parameter \emph{oppositely}, which is the most common way to
mis-set an adaptive run: raising it makes D\"orfler marking more aggressive and
fraction-of-maximum marking less so.

On the curved path there is more to it than refining. Refining $\Omega_h$
without doing anything else makes the gap to $\Gamma$ shrink more slowly than
the local element size, so the ratio the transfer analysis is stated in
\emph{grows} with every cycle and the method silently leaves the regime it is
analysed in. {\meq} therefore maintains a companion mesh, re-marks which
background elements are inside $\Gamma$ after each refinement, and rebuilds the
transfer paths --- so the domain grows as it refines, and the run reports how
many paths had to be widened.

Each cycle after the first warm-starts from the previous one, interpolated onto
the refined mesh. That costs nothing and saves a substantial fraction of the
Newton work over a run, and it is measurably neutral on the answer.

\begin{quote}
	\textbf{An error target chosen against one version of the estimator is not
	comparable against another.} The estimator is built on the enriched
	potential $\psi^{\star}$, which is a different and smaller quantity than the
	raw potential on the same mesh. If you inherit a target from somewhere,
	re-establish it before trusting it.
\end{quote}

\section{Output}
\label{sec:output}

Every run writes the same equilibrium \textbf{three times}. That is not
redundancy: no single format is simultaneously exact, portable and convenient,
and each of the three gives up a different one of those.

\begin{description}

\item[The exact format.] Three files whose names begin with the output
	stem, read by GLVis and by {\meq} itself:
	\texttt{.mesh}, \texttt{\_psi.gf} and \texttt{\_grad\_psi.gf}. They
	hold every polynomial coefficient of the finite element solution, at full
	precision. This is the restart format and the only one {\meq} can resume a
	run from. Beside them, \texttt{\_psistar.gf} carries the post-processed
	potential, which lives one degree up.

\item[The picture.] A directory named for the stem, holding a
	\texttt{.pvd} index and one subdirectory of VTK files, read by ParaView
	and VisIt.
	Lagrange cells at the degree of the field, with the mesh boundary bent onto
	the true $\Gamma$.

\item[The interchange format.] A \texttt{.nc} file, read by anything that
	reads NetCDF. The potential and the magnetic field sampled on a uniform $(R, Z)$
	grid. Lossy, and portable.

\end{description}

\subsection{Which potential each file carries}

Everything {\meq} draws or exports carries $\psi^{\star}$, the post-processed
potential, which converges a full order faster than the solved $\psi_h$ at
essentially no cost --- there is no reason to report the worse field. The one
exception is the restart file, which keeps $\psi_h$ deliberately, because it is
read back into a space of degree $k$ and $\psi^{\star}$ would not fit. Making
the two consistent would break restart.

The consequence for anyone with old output is that a gridded file differenced
against one written before this changed measures the post-processing rather than
the physics: same name, same units, same layout, different meaning. The file
records which potential it holds, and the \emph{absence} of that record is what
identifies an older one.

\subsection{The high-order picture}

The VTK output is written at the polynomial degree of the field, and this is the
one thing in the table that can be silently wrong. VTK's native cells are
linear, so the default path draws a cubic solution as though it were linear ---
which produces a picture that looks like a coarse mesh rather than like a bug,
and which nobody queries. {\meq} writes Lagrange cells instead, and its test
suite asserts the point count against the vertex count for exactly that reason.

On the curved path the same file also has its boundary bent onto the true
$\Gamma$: a curvature is installed on the mesh and each boundary face is moved
outward. Because the cells were already high order this needed nothing further
from the format; the two features composed. The gap can exceed an element's own
size, so some faces genuinely cannot reach, and the run reports the fraction
that did.

An adaptive run writes a second collection with one frame per cycle, which
ParaView scrubs through as a time series. It is separate from the answer because
the answer has its boundary bent onto $\Gamma$, and doing that mid-loop would
hand the next refinement a geometry the estimator never saw.

\subsection{The band, in the gridded file}
\label{sec:band}

The gridded file samples a uniform $(R,Z)$ grid, which on the curved path
includes nodes lying in the band --- inside the plasma, outside the mesh. Those
nodes carry real data and the file marks them.

The obvious way to fill them is to continue each field's own polynomial past its
element, and it is bounded by nothing: on a shaped test case that put a
noticeable number of nodes at \emph{positive} $\psi$, where the boundary
condition makes $\psi$ exactly zero on $\Gamma$ and strictly negative inside,
and the $\psi = 0$ contours in the band were visibly wrong. It is worth
recording that blending the extrapolation towards the known zero on $\Gamma$
does \emph{not} fix it, because it looks as though it should: scaling a positive
value down never changes its sign. The error was in \emph{where} the field was
evaluated, not in how it was weighted.

What {\meq} does instead is take a Taylor step from the node's foot on
$\Gamma_h$ using the solved flux, $\psi(p) \approx \psi(x_0) + r_0\, q(x_0)
\cdot (p - x_0)$ --- so nothing is ever evaluated outside an element, and the
step is taken in a quantity solved for at the potential's own order. The
magnetic field gets the same treatment using its own gradient. That is a full
order better than reading the value at the foot, which is what the field used to
get, and it is still not the potential's order: there is no solved variable for
the gradient of the flux, so this is second order at every degree. The gap is
structural rather than a shortcut.

\begin{quote}
	\textbf{The file says which nodes those are, and a consumer should use it.}
	Beside the mask saying which nodes are inside the plasma there is a second
	mask saying which were reached by extension. Drop those nodes before
	computing an error norm or differencing two runs: they are limited by the
	extension's order rather than by the discretisation.
\end{quote}

\subsection{What is not in any of them}

None of the three carries the flux surfaces. {\meq} can locate the magnetic
axis, trace the surfaces, take flux-surface averages over them and fit the whole
family as a smooth map from a disc --- but all of that is reached through the
C++ library, and none of it is written to a file or driven from a configuration
key. A consumer that wants $V'$, $\langle R^{-2}\rangle$ or a safety factor
links against {\meq}. That is \chref{ch:embedding}.

\section{When the solve does not converge}
\label{sec:notconverge}

This section is the one to reach for on an exit code~2.

\subsection{Read the residual history first}

It is printed on every run. Three shapes, and they mean different things.

\emph{A history that grinds down linearly}, at a roughly constant contraction
per step, means the Jacobian disagrees with the residual. If the source is one
of {\meq}'s own, that is a bug to report; if it is yours, the derivative is
wrong. Working Newton squares the residual roughly every step once it gets
going, reaching machine precision from a cold start in a handful of iterations.

\emph{A history that wanders} --- non-monotone, moving between a few orders of
magnitude and not settling --- is undamped Newton outside its basin. It is not
divergence and it is not a stall, and it is what an under-resolved mesh
produces. It is curable, and \S\ref{sec:ladder} says how.

\emph{A history that stops at zero iterations} with an identically zero residual
means the trivial branch. Several standard benchmark sources vanish at
$\psi = 0$, so with homogeneous Dirichlet data $\psi \equiv 0$ genuinely solves
the problem and Newton, which starts from the Dirichlet data, lands on it and
stops. The cure is an initial guess that puts $\psi$ away from zero in the
interior; the published algorithms for this equation open by demanding a
non-trivial initial guess for exactly this reason.

\subsection{Refine before reaching for a solver}
\label{sec:ladder}

The single most useful thing measured about {\meq}'s hard cases is that most of
them were not hard. Of four standard benchmark sources chosen for stiffness,
three converge in a handful of undamped Newton steps on a mesh that resolves
their features, and \emph{both} refinement paths cure them independently ---
halving $h$ works, and so does raising $k$ on the same mesh. A difficulty
measured at one resolution is not a property of the problem.

So the order to try things in is: resolve the feature, raise the degree, then
change the solver. Note also that the mesh has to resolve the source's own
structure and not merely look fine: a source with a ridge narrower than a cell
is a source the discretisation cannot see, and marking a generous band around
the feature and refining it once reaches a solvable configuration on fewer
degrees of freedom than uniform refinement needs.

\subsection{The globalisation ladder}

If refinement is not available --- and on the first solve of an adaptive run it
is not, because there is no estimator yet --- there is a ladder.

The default is plain undamped Newton, and every convergence rate {\meq} claims
was measured with it. The fallback the program uses on \emph{observed} failure
is a handoff: Anderson-accelerated Picard walks the iterate into Newton's basin
and plain Newton takes it from there and supplies the quadratic endgame Picard
structurally cannot. It reaches cases nothing else in the solver touches at that
resolution, and it agrees with Picard's own answer, which is what makes it a
handoff rather than a change of problem. Plain and accelerated Picard are
available on their own, as is the SUNDIALS nonlinear solver with and without its
line search.

Three things about the ladder are worth knowing before adjusting it.

\emph{Picard's job is not to solve the problem.} Stage one stopping short of its
own tolerance is an expected outcome and the handoff swallows it deliberately.
This is a globalisation, not a two-solver pipeline.

\emph{Do not replace Picard's tolerance with an iteration budget.} The handoff
is not monotone in Picard effort --- budgets both smaller and larger than an
intermediate one converge where the intermediate one fails --- so a budget tuned
on one mesh will betray you on the next. Picard's own tolerance is the trigger
that works wherever it is reached.

\emph{The fallback is reactive, never predictive.} Two candidate cheap
predictors of failure have been measured and both killed, and the more obvious
of the two --- the first Newton step making the residual worse --- turned out to
be \emph{anti}-correlated: the largest first-step blow-up measured converges,
and the cases where the residual came down at step one both fail. Nothing may be
inferred from the source about which solver to run.

\subsection{Why not make the fallback the default}

Because it silently changes which discrete solution you get, and that settles it
on its own. Coarse discretisations of these sources carry \emph{more than one}
solution: the same case solved to a tight tolerance by three different routes
gives three converged answers differing by a noticeable fraction in the peak
flux. These are not tolerance artefacts --- tightening the tolerance leaves them
bit-identical. So changing the default globalisation is not a performance
decision; it changes the equilibrium {\meq} reports on an under-resolved mesh,
which is precisely the mesh an adaptive run starts from.

The corollary for anyone comparing two solve routes is: do it on a mesh that
resolves the problem. A disagreement on a coarse one is this, not a defect.

\section{Options settled by measurement}

Several choices in {\meq} are exposed because the right answer depends on the
problem and the machine. In each case the default was chosen by measurement
here, and in each case the standing advice is the same: \emph{measure it on your
own problem before departing from it.}

\begin{description}

\item[The trace solver.] The globally coupled trace system is solved by a sparse
	direct method, and there are up to three available depending on what MFEM
	was built with. They agree with one another to round-off through the whole
	solver, so choosing between them is never a numerical decision --- it is a
	licensing and a performance one. The default is the one present in every
	build.

\item[Threaded assembly.] The element-local half of the assembly can be
	threaded, and it reproduces the serial result bit for bit, which is
	assertable rather than merely tolerable because element-local arithmetic
	reassociates nothing. It is nevertheless \emph{off} by default, and the
	reason is instructive: the library forks a thread team per call, so a caller
	that assembles once amortises it and a caller that assembles hundreds of
	times inside an inner iteration pays it every time. Mesh size does not
	separate the two cases and the solver cannot know which caller it has. Take
	it for a large mesh assembled a few times; leave it for a small one
	assembled often.

\item[The nonlinear ordering.] Whether the system is linearised before it is
	hybridized or after. The default linearises first, in which case every
	element-local operation is a linear solve; the alternative condenses first
	and puts a small nonlinear solve inside every element for every residual
	evaluation. The first is several times cheaper per iteration and its
	per-element work is uniform, which makes it cheap even when it is failing.
	The second is more robust on stiff, under-resolved meshes, and is kept as
	the backup for exactly that. \emph{Do not expect an ordering to fix a stiff
	source} --- that prediction was made here, tested and falsified.

\item[Threads in the numerical library.] If MFEM is linked against a threaded
	BLAS, the number of threads it is allowed matters enormously and not in the
	direction one expects. The element-local dense factorisations are small, and
	handing a small dense factorisation to a threaded library costs a fork and a
	barrier per call for no work; at high polynomial degree, where each such
	block crosses the library's own threading threshold, the cost is
	catastrophic and it looks exactly like a solver problem. {\meq}'s own tests
	all pin the thread count to one. \textbf{A production run at high degree
	with threading left on would be unusable.}

\end{description}

Two remarks that apply to all of them. The largest performance lever in {\meq}
is none of these --- it is resolution, and after that the globalisation, and on
a hard case the nonlinear iteration dominates everything in this list. And a
timing taken on one machine is a measurement about that machine: this is
especially true of anything involving a graphics processor, where a consumer
card runs double precision at a small fraction of a datacentre part's rate and
can invert the conclusion outright.

\section{Tools}

The source tree carries a small plotting script that reads the gridded output,
and a short guide saying which of the three formats goes with which reader and
why they are not interchangeable. Neither is part of the solver, and the guide
is the thing to read first if a picture looks wrong.
